\SourceLeaf{410}%
\par
Man kann dieses auch durch folgende sechs Gleichungen ausdrücken:
\[
\begin{aligned}
\frac{\sin A}{\sin a}=\frac{\sin B}{\sin b}=\frac{\sin C}{\sin c}&=-\frac{\cos\tfrac12(A+B+C)}{2\sin\tfrac12a.\sin\tfrac12b.\sin\tfrac12c}\\
=\frac{\cos\tfrac12(-A+B+C)}{2\sin\tfrac12a.\cos\tfrac12b.\cos\tfrac12c}&=\frac{\cos\tfrac12(A-B+C)}{2\cos\tfrac12a.\sin\tfrac12b.\cos\tfrac12c}=\frac{\cos\tfrac12(A+B-C)}{2\cos\tfrac12a.\cos\tfrac12b.\sin\tfrac12c}
\end{aligned}
\]
oder durch folgende
\[
\begin{aligned}
\sin A\sin b\sin c&=\sin B\sin a\sin c=\sin C\sin a\sin b\\
&=-4\cos\tfrac12(+A+B+C).\cos\tfrac12a.\cos\tfrac12b.\cos\tfrac12c\\
&=4\cos\tfrac12(-A+B+C).\cos\tfrac12a.\sin\tfrac12b.\sin\tfrac12c\\
&=4\cos\tfrac12(+A-B+C).\sin\tfrac12a.\cos\tfrac12b.\sin\tfrac12c\\
&=4\cos\tfrac12(+A+B-C).\sin\tfrac12a.\sin\tfrac12b.\cos\tfrac12c
\end{aligned}
\]
Diese Grösse bedeutet den sechsfachen Inhalt der Pyramide, deren Ecken die drei Winkelpunkte des sphärischen Dreiecks und der Mittelpunkt der Kugel bilden, Halbmesser der Kugel \(=1\) gesetzt. Ferner ist diese Grösse
\[
=4\operatorname{cotang}r.\sin\tfrac12a.\sin\tfrac12b.\sin\tfrac12c
\]
wo \(r\) den sphärischen Halbmesser des um das Dreieck beschriebenen Kreises bedeutet. Auch ist dieselbe
\[
=4\cos r.\sin\alpha.\sin\tfrac12b.\sin\tfrac12c
\]
wo [\(2\alpha,2\beta,2\gamma\) die Winkel bedeuten, welche zwischen je zwei der nach den Eckpunkten \(A,B,C\) gezogenen sphärischen Halbmesser und gegenüber den Seiten \(a,b,c\) liegen, oder] \(\alpha,\beta,\gamma\) die Winkel des ebenen Dreiecks \(ABC\) sind, weil \(2\sin\tfrac12a,\ 2\sin\tfrac12b,\ 2\sin\tfrac12c\) dessen Seiten, mithin \(4\sin\alpha.\sin\tfrac12b.\sin\tfrac12c\) dessen doppelter Inhalt; während dasselbe zugleich als Grundfläche obiger Pyramide mit Höhe \(\cos r\) betrachtet werden kann, woraus die Richtigkeit von selbst erhellt. [Aus der obigen sechsfachen Gleichung leitet \textsc{Gauss} an einer andern Stelle die von ihm so vielfach angewandten Formeln her:]
\[
\begin{aligned}
\cos\tfrac12A.\sin\tfrac12(b-c)&=\sin\tfrac12(B-C).\sin\tfrac12a\\
\sin\tfrac12A.\sin\tfrac12(b+c)&=\cos\tfrac12(B-C).\sin\tfrac12a\\
\cos\tfrac12A.\cos\tfrac12(b-c)&=\sin\tfrac12(B+C).\cos\tfrac12a\\
\sin\tfrac12A.\cos\tfrac12(b+c)&=\cos\tfrac12(B+C).\cos\tfrac12a
\end{aligned}
\]
